\begin{figure}[ht]
  \centering
  \begin{tikzpicture}
    \begin{axis}[paperaxis,
      xlabel={Année}, ylabel={Écart France moins comparables (variation du log capital)},
      xtick={2014,2015,2016,2017,2018,2019,2020,2021,2022},
      xticklabel style={/pgf/number format/1000 sep=},
    ]
      \draw[zeroline, line width=0.6pt] (axis cs:2013.5,0) -- (axis cs:2022.5,0);
      % La France, en bleu, avec intervalle a 95 pour cent.
      \addplot[franceblue, mark=*, mark options={scale=0.8}, error bars/.cd,
        y dir=both, y explicit, error bar style={franceblue}]
        table[col sep=comma, x=year, y=coef, y error expr=1.96*\thisrow{std_error}]
        {figures/entree_capital.csv};
      % Reforme de 2018 : marqueur vert, la couleur de la baisse d'impot.
      \draw[baisse2018, dashed, line width=1pt]
        ({axis cs:2017.5,0}|-{rel axis cs:0,0}) -- ({axis cs:2017.5,0}|-{rel axis cs:0,1});
    \end{axis}
  \end{tikzpicture}
  \caption{Étude d'événement de l'entrée de capital autour de la réforme de 2018}
  \label{fig:entree}
  \fignote{Coefficients $\tau_j$ de l'équation~\eqref{eq:event} pour la variation du logarithme du
  capital souscrit, écart de la France aux cinq pays comparables l'année $j$, année de référence 2017 omise.
  Barres d'erreur à quatre-vingt-quinze pour cent sur l'erreur-type regroupée par pays, indicative.
  La ligne verte verticale marque la réforme de 2018. Un profil plat avant et après la réforme
  indique l'absence d'entrée de capital neuf.}
\end{figure}
